%!TeX root=../persuasiontop.tex
\chapter[Chapter \thechapter]{}

\begin{figure}[t!]
\centering
\includegraphics[width=\linewidth]{nouveau24}
\end{figure}

\lettrine[lines=4,lraise=0.3]{W}{ho} can be in doubt of what followed? When any two young people take it into their heads to marry, they are pretty sure by perseverance to carry their point, be they ever so poor, or ever so imprudent, or ever so little likely to be necessary to each other's ultimate comfort. This may be bad morality to conclude with, but I believe it to be truth; and if such parties succeed, how should a Captain Wentworth and an Anne Elliot, with the advantage of maturity of mind, consciousness of right, and one independent fortune between them, fail of bearing down every opposition? They might in fact, have borne down a great deal more than they met with, for there was little to distress them beyond the want of graciousness and warmth. Sir Walter made no objection, and Elizabeth did nothing worse than look cold and unconcerned. Captain Wentworth, with five-and-twenty thousand pounds, and as high in his profession as merit and activity could place him, was no longer nobody. He was now esteemed quite worthy to address the daughter of a foolish, spendthrift baronet, who had not had principle or sense enough to maintain himself in the situation in which Providence had placed him, and who could give his daughter at present but a small part of the share of ten thousand pounds which must be hers hereafter.

Sir Walter, indeed, though he had no affection for Anne, and no vanity flattered, to make him really happy on the occasion, was very far from thinking it a bad match for her. On the contrary, when he saw more of Captain Wentworth, saw him repeatedly by daylight, and eyed him well, he was very much struck by his personal claims, and felt that his superiority of appearance might be not unfairly balanced against her superiority of rank; and all this, assisted by his well-sounding name, enabled Sir Walter at last to prepare his pen, with a very good grace, for the insertion of the marriage in the volume of honour.

The only one among them, whose opposition of feeling could excite any serious anxiety was Lady Russell. Anne knew that Lady Russell must be suffering some pain in understanding and relinquishing Mr Elliot, and be making some struggles to become truly acquainted with, and do justice to Captain Wentworth. This however was what Lady Russell had now to do. She must learn to feel that she had been mistaken with regard to both; that she had been unfairly influenced by appearances in each; that because Captain Wentworth's manners had not suited her own ideas, she had been too quick in suspecting them to indicate a character of dangerous impetuosity; and that because Mr Elliot's manners had precisely pleased her in their propriety and correctness, their general politeness and suavity, she had been too quick in receiving them as the certain result of the most correct opinions and well-regulated mind. There was nothing less for Lady Russell to do, than to admit that she had been pretty completely wrong, and to take up a new set of opinions and of hopes.

There is a quickness of perception in some, a nicety in the discernment of character, a natural penetration, in short, which no experience in others can equal, and Lady Russell had been less gifted in this part of understanding than her young friend. But she was a very good woman, and if her second object was to be sensible and well-judging, her first was to see Anne happy. She loved Anne better than she loved her own abilities; and when the awkwardness of the beginning was over, found little hardship in attaching herself as a mother to the man who was securing the happiness of her other child.

Of all the family, Mary was probably the one most immediately gratified by the circumstance. It was creditable to have a sister married, and she might flatter herself with having been greatly instrumental to the connexion, by keeping Anne with her in the autumn; and as her own sister must be better than her husband's sisters, it was very agreeable that Captain Wentworth should be a richer man than either Captain Benwick or Charles Hayter. She had something to suffer, perhaps, when they came into contact again, in seeing Anne restored to the rights of seniority, and the mistress of a very pretty landaulette; but she had a future to look forward to, of powerful consolation. Anne had no Uppercross Hall before her, no landed estate, no headship of a family; and if they could but keep Captain Wentworth from being made a baronet, she would not change situations with Anne.

It would be well for the eldest sister if she were equally satisfied with her situation, for a change is not very probable there. She had soon the mortification of seeing Mr Elliot withdraw, and no one of proper condition has since presented himself to raise even the unfounded hopes which sunk with him.

The news of his cousin Anne's engagement burst on Mr Elliot most unexpectedly. It deranged his best plan of domestic happiness, his best hope of keeping Sir Walter single by the watchfulness which a son-in-law's rights would have given. But, though discomfited and disappointed, he could still do something for his own interest and his own enjoyment. He soon quitted Bath; and on Mrs Clay's quitting it soon afterwards, and being next heard of as established under his protection in London, it was evident how double a game he had been playing, and how determined he was to save himself from being cut out by one artful woman, at least.

Mrs Clay's affections had overpowered her interest, and she had sacrificed, for the young man's sake, the possibility of scheming longer for Sir Walter. She has abilities, however, as well as affections; and it is now a doubtful point whether his cunning, or hers, may finally carry the day; whether, after preventing her from being the wife of Sir Walter, he may not be wheedled and caressed at last into making her the wife of Sir William.

It cannot be doubted that Sir Walter and Elizabeth were shocked and mortified by the loss of their companion, and the discovery of their deception in her. They had their great cousins, to be sure, to resort to for comfort; but they must long feel that to flatter and follow others, without being flattered and followed in turn, is but a state of half enjoyment.

Anne, satisfied at a very early period of Lady Russell's meaning to love Captain Wentworth as she ought, had no other alloy to the happiness of her prospects than what arose from the consciousness of having no relations to bestow on him which a man of sense could value. There she felt her own inferiority very keenly. The disproportion in their fortune was nothing; it did not give her a moment's regret; but to have no family to receive and estimate him properly, nothing of respectability, of harmony, of good will to offer in return for all the worth and all the prompt welcome which met her in his brothers and sisters, was a source of as lively pain as her mind could well be sensible of under circumstances of otherwise strong felicity. She had but two friends in the world to add to his list, Lady Russell and Mrs Smith. To those, however, he was very well disposed to attach himself. Lady Russell, in spite of all her former transgressions, he could now value from his heart. While he was not obliged to say that he believed her to have been right in originally dividing them, he was ready to say almost everything else in her favour, and as for Mrs Smith, she had claims of various kinds to recommend her quickly and permanently.

Her recent good offices by Anne had been enough in themselves, and their marriage, instead of depriving her of one friend, secured her two. She was their earliest visitor in their settled life; and Captain Wentworth, by putting her in the way of recovering her husband's property in the West Indies, by writing for her, acting for her, and seeing her through all the petty difficulties of the case with the activity and exertion of a fearless man and a determined friend, fully requited the services which she had rendered, or ever meant to render, to his wife.

Mrs Smith's enjoyments were not spoiled by this improvement of income, with some improvement of health, and the acquisition of such friends to be often with, for her cheerfulness and mental alacrity did not fail her; and while these prime supplies of good remained, she might have bid defiance even to greater accessions of worldly prosperity. She might have been absolutely rich and perfectly healthy, and yet be happy. Her spring of felicity was in the glow of her spirits, as her friend Anne's was in the warmth of her heart. Anne was tenderness itself, and she had the full worth of it in Captain Wentworth's affection. His profession was all that could ever make her friends wish that tenderness less, the dread of a future war all that could dim her sunshine. She gloried in being a sailor's wife, but she must pay the tax of quick alarm for belonging to that profession which is, if possible, more distinguished in its domestic virtues than in its national importance.

\begin{figure}[b!]
\centering
\includegraphics[width=.4\linewidth]{tailpiece}
\end{figure}